% Generated draft: requires mathematical and source review.
\documentclass[11pt]{article}
\usepackage[margin=25mm]{geometry}
\usepackage{fontspec}
\setmainfont{Latin Modern Roman}
\usepackage{amsmath,amssymb,mathtools,unicode-math}
\setmathfont{Latin Modern Math}
\usepackage{enumitem,needspace,xurl,hyperref}
\setlength{\parindent}{0pt}
\setlength{\parskip}{5pt}
\setlength{\emergencystretch}{3em}
\newcommand{\N}{\mathbb{N}}
\newcommand{\Z}{\mathbb{Z}}
\newcommand{\Q}{\mathbb{Q}}
\newcommand{\R}{\mathbb{R}}
\newcommand{\C}{\mathbb{C}}
\newcommand{\F}{\mathbb{F}}
\newcommand{\A}{\mathbb{A}}
\newcommand{\PP}{\mathbb{P}}
\newcommand{\E}{\mathbb{E}}
\newcommand{\eps}{\varepsilon}
\newcommand{\dd}{\,\mathrm{d}}
\begin{document}
\section*{UnsolvedMath ID: 30000808}
\textbf{Source problem code:} catalogue entry
\subsection{Definitions and mathematical statement}
Let $F$ be an algebraically closed field of characteristic $\ell$; the record does not expressly exclude characteristic zero. For a finite group $K$, let $Z(K)$ denote its center. Say that $K/Z(K)$ is almost simple if, for some nonabelian finite simple group $S$, \[S\leq K/Z(K)\leq\operatorname{Aut}(S)\] after identifying $S$ with its inner automorphism group. Classify all triples $(K,V,H)$ with the following properties: $K$ is finite, $K/Z(K)$ is almost simple, $V$ is an $FK$-module with $\dim_FV>1$, $H$ is a proper subgroup of $K$, and $V\!\downarrow_H$ is irreducible. Here $V\!\downarrow_H$ is the same vector space with the action restricted to $H$; irreducibility means that its only $H$-invariant $F$-linear subspaces are $0$ and $V$. The finiteness of $K$, explicit in the original statement but omitted in the clean/current wording, is retained. No faithfulness, perfectness or quasisimplicity of $K$, maximality of $H$, or cross-characteristic assumption is added. In particular the question concerns all proper subgroups and is not restricted to maximal ones. A classification may be expressed in terms of group isomorphisms carrying $H$ to the corresponding subgroup and intertwining the modules, but the source does not specify its equivalence convention. The supplied 2026-08-21 triage labels the problem partially solved and reports substantial completed cases, including general linear groups in cross characteristic and symmetric or alternating groups, while describing the full almost-simple classification across characteristics as active.
\subsection{Short English statement}
Classify the finite groups with almost-simple central quotient, their modules of dimension greater than one, and proper subgroups on which those modules remain irreducible. The original's finite-group hypothesis is retained, without adding faithfulness or a maximal-subgroup restriction.
\subsection{Sources}
\begin{itemize}\item Local UnsolvedMath registry, permanent ID 30000808, source problem catalogue entry, \emph{Irreducible Restrictions of Almost-Simple Group Modules}, proposed year 2007; no individual proposer supplied.\item Supplied citation: \emph{Permutation Groups} (2008), Oberwolfach Reports, DOI 10.4171/owr/2007/37, pp.\ 2173--2236: \url{https://doi.org/10.4171/owr/2007/37}. Supplied literature label: original Oberwolfach report.\item Supplied literature label: Kleshchev--Tiep, general linear groups, \url{https://doi.org/10.1353/ajm.0.0108}.\item Supplied literature label: Kleshchev--Morotti--Tiep, small characteristics, \url{https://arxiv.org/abs/1903.09854}.\end{itemize}
\end{document}
